\section{实验与证据}

\begin{frame}{负结果同样决定工程选型}
  \begin{columns}[c,onlytextwidth]
    \begin{column}{0.62\textwidth}
      \centering
      \includegraphics[width=\linewidth,height=0.64\textheight,keepaspectratio]
        {terrain_following/terrain_clearance_rmse_pid_mpc.png}
    \end{column}
    \begin{column}{0.34\textwidth}
      \begin{block}{近底地形跟随}
        三档地形、每档三种子均无安全违规；调优 PID 内层不弱于 MPC。
      \end{block}
      \begin{block}{路径跟踪}
        MPC 在长/短 S 弯和发卡路径优于或持平基线；直角折弯上 LOS 前瞻更优。
      \end{block}
      \begin{alertblock}{研究结论}
        不存在“一个控制器覆盖所有工况”的证据。
      \end{alertblock}
    \end{column}
  \end{columns}
\end{frame}

\begin{frame}{分层证据：结论强度不超过验证层级}
  \centering
  \includegraphics[width=0.88\linewidth,height=0.75\textheight,keepaspectratio]
    {architecture/auv_system_verification_deployment_ladder.png}

  \FrameConclusion{本文已闭合算法、数字孪生和部署接口证据；实物标定与海试仍是独立门槛。}
\end{frame}

\begin{frame}{从先验畸变失败到六自由度闭环恢复}
  \begin{columns}[c,onlytextwidth]
    \begin{column}{0.70\textwidth}
      \centering
      \includegraphics[width=\linewidth,height=0.62\textheight,keepaspectratio]
        {cable_acceptance/pvs_closedloop_recovery_prior_alignment.png}
    \end{column}
    \begin{column}{0.27\textwidth}
      \small
      \begin{alertblock}{先暴露失败}
        中、重档开环回放各 \(3/3\) 不通过，偏差传导至验收。
      \end{alertblock}
      \begin{block}{物理前提修复}
        磁观测被真实接受后启用在线修正；阈值不变。
      \end{block}
      \begin{block}{闭环恢复}
        中、重档各 \(3/3\) 通过；重档约 12 s 进入 \(\pm3.4\) m 廊道。
      \end{block}
      \scriptsize
      \BoundaryTag{边界}\quad 数字孪生、静态畸变、满足磁观测前提。
    \end{column}
  \end{columns}
\end{frame}
